\chapter{Smoking (in pregnancy)}

\medskip
\noindent\textit{\textbf{Under review.} This chapter is an AI-assisted educational draft; it carries no source citations and is not a substitute for professional medical advice.}

\section*{Definition and Overview}
Smoking (in pregnancy) relates to women's reproductive and sexual health, encompassing menstruation, fertility, pregnancy, childbirth, and menopause, as well as conditions affecting the female reproductive system. These events and conditions reflect normal physiology as well as pathology, and management requires attention to both biomedical and psychosocial factors. Care is provided by GPs, midwives, obstetricians, and gynaecologists.

\section*{Epidemiology}
Menstrual disorders affect up to 25\% of women of reproductive age. Endometriosis affects approximately 1 in 10 women. Polycystic ovary syndrome (PCOS) affects 8--13\% of reproductive-aged women. Approximately 10--15\% of couples experience infertility. Menopause typically occurs at age 50--52. Cervical cancer incidence has fallen significantly with screening and HPV vaccination programmes.

\section*{Aetiology and Pathophysiology}
\begin{itemize}
    \item \textbf{Hormonal changes:} fluctuations in oestrogen, progesterone, FSH, and LH across the menstrual cycle, pregnancy, and menopause
    \item \textbf{Structural abnormalities:} fibroids, polyps, endometriosis, ovarian cysts, congenital uterine anomalies
    \item \textbf{Infection:} sexually transmitted infections, pelvic inflammatory disease, vaginal candidiasis
    \item \textbf{Autoimmune/endocrine:} thyroid disorders, premature ovarian insufficiency
    \item \textbf{Genetic:} hereditary cancer syndromes (BRCA1/2, Lynch syndrome)
    \item \textbf{Lifestyle factors:} stress, weight extremes, excessive exercise, smoking
\end{itemize}

\section*{Clinical Features}
\begin{itemize}
    \item Changes in menstrual cycle: heavy, irregular, or absent periods
    \item Dysmenorrhoea (painful periods) or dyspareunia (painful intercourse)
    \item Intermenstrual or postcoital bleeding
    \item Postmenopausal bleeding -- always requires investigation
    \item Abnormal vaginal discharge (colour, odour, or itch)
    \item Pelvic pain, pressure, or a palpable mass
    \item Menopausal symptoms: hot flushes, night sweats, mood changes, vaginal dryness
    \item Subfertility or difficulty conceiving
\end{itemize}

\section*{Differential Diagnosis}
\begin{itemize}
    \item Pregnancy-related: ectopic pregnancy, miscarriage, molar pregnancy
    \item Structural: fibroids, endometrial polyps, ovarian cysts, endometriosis
    \item Hormonal: PCOS, thyroid dysfunction, hyperprolactinaemia
    \item Infective: pelvic inflammatory disease, cervicitis, vaginitis
    \item Neoplastic: cervical, endometrial, or ovarian cancer
    \item Functional: dysfunctional uterine bleeding, primary dysmenorrhoea
\end{itemize}

\section*{When to Seek Medical Care}
See a GP for menstrual changes, persistent pelvic pain, unusual bleeding, or unusual discharge.

\textbf{Call 000} for:
\begin{itemize}
    \item Heavy vaginal bleeding soaking pads rapidly, especially in pregnancy
    \item Severe abdominal or pelvic pain
    \item Fever with pelvic pain or discharge (possible pelvic infection or sepsis)
    \item Shoulder tip pain with vaginal bleeding and positive pregnancy test (suspected ectopic pregnancy)
    \item Bleeding or severe pain after miscarriage, abortion, or childbirth
\end{itemize}

\section*{Diagnosis and Investigations}
\begin{itemize}
    \item \textbf{History:} menstrual, obstetric, gynaecological, sexual, and contraceptive history
    \item \textbf{Examination:} abdominal and pelvic (bimanual and speculum) examination
    \item \textbf{Microbiology:} swabs for chlamydia, gonorrhoea, and other STIs; wet mount for candidiasis
    \item \textbf{Blood tests:} beta-hCG (pregnancy test), full blood count, iron studies, thyroid function, hormone profile (FSH, LH, oestradiol, prolactin, testosterone)
    \item \textbf{Imaging:} pelvic ultrasound (transvaginal or transabdominal); saline infusion sonography; MRI for complex cases
    \item \textbf{Histology:} endometrial biopsy, hysteroscopy with biopsy, colposcopy with cervical biopsy
\end{itemize}

\section*{Management}
\begin{itemize}
    \item \textbf{Lifestyle:} weight management, exercise, stress reduction
    \item \textbf{Hormonal therapy:} combined oral contraceptive pill, progesterone-only pill, Mirena IUD, hormone replacement therapy (HRT) for menopause
    \item \textbf{Non-hormonal medications:} NSAIDs for dysmenorrhoea, antifibrinolytics (tranexamic acid) for heavy bleeding
    \item \textbf{Antimicrobials:} antibiotics for STIs and pelvic inflammatory disease; antifungals for candidiasis
    \item \textbf{Surgery:} myomectomy or hysterectomy for fibroids; endometrial ablation for heavy menstrual bleeding; laparoscopic surgery for endometriosis
    \item \textbf{Fertility management:} ovulation induction, intrauterine insemination, IVF
    \item \textbf{Pelvic floor physiotherapy:} for prolapse or incontinence
\end{itemize}

\section*{Complications}
\begin{itemize}
    \item Iron-deficiency anaemia from heavy menstrual bleeding
    \item Infertility or subfertility
    \item Chronic pelvic pain
    \item Ectopic pregnancy and rupture (life-threatening)
    \item Sepsis from untreated pelvic infection
    \item Endometrial hyperplasia and progression to endometrial cancer (with unopposed oestrogen)
    \item Osteoporosis and cardiovascular disease (postmenopausal)
    \item Psychological impact of menopausal symptoms or infertility
\end{itemize}

\section*{Prognosis and Outcomes}
Most menstrual disorders are well controlled with medical or surgical treatment. Endometriosis may require long-term management. PCOS is a chronic condition managed with lifestyle measures and symptom-targeted therapy. Menopausal symptoms typically resolve over 2--5 years, though bone and cardiovascular risks persist. Fertility outcomes depend on the underlying cause and treatment. Early detection and treatment of reproductive cancers improves survival significantly.

\section*{Prevention and Self-Care}
\begin{itemize}
    \item Participate in cervical screening (HPV test every 5 years for women aged 25--74)
    \item HPV vaccination (ideally before sexual debut)
    \item Maintain a healthy weight and exercise regularly
    \item Use contraception to plan pregnancy and prevent STIs
    \item Take folic acid supplementation before and during early pregnancy
    \item Avoid smoking and limit alcohol, especially in pregnancy
    \item Seek pre-pregnancy care if planning a pregnancy
    \item Discuss HRT risks and benefits with a doctor during menopause transition
\end{itemize}
